%=====================================================%
%           CHAPTER 2: LITERATURE REVIEW
%=====================================================%

\section{Literature Review}

Recent studies show that AI-based crop disease detection systems are becoming one of the most reliable solutions for modern farming challenges. Research highlights that traditional visual inspection by farmers is often slow, inaccurate, and highly dependent on experience. To overcome these issues, recent literature strongly supports the use of deep learning, especially Convolutional Neural Networks (CNNs), for analyzing leaf images and identifying disease patterns automatically.

Patel \& Singh (2022) proposed a CNN-based model that extracts fine-grained features such as colour variations, textures, and lesion shapes, providing significantly higher accuracy compared to classical machine learning techniques. Several researchers also emphasize the importance of image datasets, data augmentation, and training optimization methods to improve model robustness under varying lighting and environmental conditions. Studies further show that integrating AI models with cloud platforms or mobile applications allows farmers to receive instant disease predictions directly in the field. Overall, the literature agrees that AI-driven disease detection systems offer faster diagnosis, reduce crop losses, and support the development of an efficient and scalable smart agriculture platform, which closely aligns with the objectives of the proposed project \cite{ref1}.

Agriculture in India faces several challenges such as outdated farming methods, climate uncertainty, and limited access to modern technology. Recent studies have shown that Artificial Intelligence and Machine Learning can improve agricultural productivity through crop prediction, disease detection, and yield forecasting. Researchers such as Rahman, Shamrat, Tasnim, Roy, and Hossain (2019) highlighted the effectiveness of supervised Machine Learning algorithms in intelligent prediction systems and decision-making models \cite{ref2}.

Many existing agricultural systems provide separate services such as weather forecasting, crop recommendation, or online trading, but they often lack an integrated and user-friendly platform for farmers. To overcome these limitations, the proposed Agricultural Assistance \& Trading Portal combines Machine Learning models, AI-powered tools, weather forecasting, agricultural news, and crop trading features into a single platform. The system helps farmers make informed decisions, improve productivity, and access modern agricultural technologies more efficiently.

Finally, the Agricultural Assistance \& Trading Portal demonstrates how technologies such as Machine Learning, Artificial Intelligence, cloud APIs, and full-stack web development can be effectively applied in the agricultural sector. The system provides intelligent decision-support tools, secure digital trading, and real-time agricultural assistance through an integrated platform. By improving accessibility to smart farming technologies and reducing the gap between farmers and consumers, the project contributes toward efficient, sustainable, and technology-driven agriculture.

Existing research in smart agriculture highlights the importance of Machine Learning and data-driven systems in supporting farmers' decision-making processes. Various studies have used parameters such as temperature, rainfall, soil type, season, and cultivation area for crop and yield prediction. Researchers have applied algorithms including Artificial Neural Networks (ANN), Regression Models, K-Nearest Neighbors (KNN), Decision Trees, Random Forests, and Stacking Regression to improve prediction accuracy and agricultural planning.

Several studies also focus on fertilizer recommendation, weather forecasting, and crop suitability analysis using soil and climate data. Researchers have emphasized the use of association-rule mining and predictive analytics to identify relationships between agricultural parameters and farming outcomes. Liu, Zhang, and Wang (2025) discussed the recent advancements and challenges in Machine Learning-based crop yield prediction systems and highlighted the growing role of intelligent agriculture technologies in improving productivity and sustainability \cite{ref3}.

These research contributions demonstrate that integrating Machine Learning models with web-based platforms can significantly enhance agricultural efficiency and farmer support systems. Inspired by these advancements, the proposed Agricultural Assistance \& Trading Portal combines crop prediction, yield estimation, weather forecasting, fertilizer recommendation, AI chatbot assistance, and online crop trading into a single integrated platform. The system aims to provide real-time recommendations, intelligent decision support, and improved accessibility to modern agricultural technologies for farmers and customers.

Recent research highlights the growing importance of Artificial Intelligence, Internet of Things (IoT), and precision agriculture in achieving sustainable and climate-smart farming. Technologies such as Machine Learning, deep learning, and remote sensing are widely used for soil analysis, irrigation optimization, crop monitoring, and disease detection. Precision farming tools including GPS-based systems and Variable Rate Technology (VRT) help improve resource utilization, reduce wastage, and increase agricultural productivity. Manjula and Djodiltachoumy (2017) proposed a crop yield prediction model that demonstrated the effectiveness of computational intelligence techniques in improving agricultural planning and productivity \cite{ref4}.

However, many existing agricultural systems focus only on individual services and lack a centralized, user-friendly platform that integrates multiple functionalities. To overcome this limitation, the proposed Agricultural Assistance \& Trading Portal combines Machine Learning models, AI chatbot support, crop prediction, weather forecasting, fertilizer recommendation, and online crop trading within a single platform. The system aims to provide intelligent, accessible, and technology-driven agricultural support that promotes sustainable farming and improves farmer productivity.

Agriculture is one of the most important sectors for economic growth and food sustainability, but it faces major challenges due to climate change, unpredictable weather conditions, and limited agricultural resources. Recent studies highlight the growing use of Artificial Intelligence (AI) and Machine Learning (ML) in improving agricultural productivity and decision-making. Techniques such as supervised learning, unsupervised learning, reinforcement learning, and Long Short-Term Memory (LSTM) models are widely used for crop yield prediction, weather forecasting, irrigation planning, and resource optimization. Karishma Kumari, Ali Mirzakhani Nafchi, Salman Mirzaee, and Ahmed Abdalla (2025) discussed the role of AI-driven technologies in achieving climate-smart and sustainable agriculture through intelligent prediction and analysis systems \cite{ref5}.

Despite these advancements, many existing agricultural solutions focus only on specific functionalities and do not provide a centralized platform for complete farm management. To address this limitation, the proposed Agricultural Assistance \& Trading Portal integrates crop prediction, fertilizer recommendation, yield estimation, weather forecasting, AI chatbot assistance, agricultural news, and online crop trading into a single platform. The system provides a multilingual, intelligent, and user-friendly solution that supports modern farming practices, improves decision-making, and promotes sustainable agriculture through technology-driven assistance.
